\documentclass[10pt]{article}

\usepackage[utf8]{inputenc}

\usepackage[T1]{fontenc}

\usepackage{amsmath}

\usepackage{amsfonts}

\usepackage{amssymb}

\usepackage[version=4]{mhchem}

\usepackage{stmaryrd}

\usepackage{graphicx}

\usepackage[export]{adjustbox}

\graphicspath{ {./images/} }

\usepackage{bbold}



\DeclareUnicodeCharacter{0131}{$\imath$}



\begin{document}

$\rightarrow M$ - векторное расслоение, ассоциированное с л и $G$-модулем $E$. Связность $\nabla$ расслоения $\pi$ определяет оператор $\nabla^{E}: \Gamma(E) \rightarrow \Gamma\left(T^{*} M \otimes E M\right)$, действующий в пространстве $\Gamma(E)$ сечений расслоения $E(M)$. Он продолжается до оператора $d^{\nabla}: \Gamma\left(E_{p}\right) \rightarrow \Gamma\left(E_{p+1}\right)$, дөйствующего в пространстве $\Gamma\left(E_{p}\right) E(M)$-значных $p$-форм по формуле $d^{\nabla}(\omega \otimes e)=d \omega(\bigotimes) e+(-1)^{p}(1) \wedge \nabla^{E} e$. Формально сопряженный к $d^{\nabla}$ оператор $\delta^{\nabla}$ на $p$-формах равен $\delta^{\nabla}=(-1)^{p+1} * d^{\nabla} *$, где $*$ - оператор Ходжа,



Связность $\nabla$ в главном $G$-расслоении наз. п о ле м Я нг а - М и л с а, если кривизна $R \nabla$, рассматриваемая как 2-форма со значениями в векторном расслоении $\because(M)=P \times$ AdGg (где $\mathfrak{g}-$ алгебра Ли группы Ли $G)$, удовлетворяөт уравнению $\delta^{\nabla} R^{\nabla}=0$.



Для римановой связности $\nabla^{g}$ риманова многообразия $(M, g)$ уравнение Янга - Миллса равносильно условию симметричности



$\left(\nabla_{X}^{g}\right.$ ric $)(Y, Z)=\left(\nabla_{Y}^{g}\right.$ ric $)(X, Z), X, Y, Z \in D(M)$ ковариантной производной тензора Риччи ric. Таким образом, примерами Я. - М. П. служат римановы связности пространств Эйнштейна и прямых произведений таких пространств. В частности, римановы связности пространств Кәлера - Эйнштөйна и кватернионных римановых пространств определяют Я.- М. П. в главных расслоениях реперов со структурными группами $U(n / 2)$ и $S p(1) \cdot S p(n / 4)$, Примерами неәйнштейновых римановых связностей, удовлетворяюших уравнению Янга - Миллса, являются римановы связности конформно плоских метрик с постоянной скалярної кривизной и непостоянной секционной кривизной. Примерами неримановых связностей, удовлетворяющих уравнениям Янга - Миллса, являются связности в нормальном расслоении вполне геодезич. подмногообразиї симметрического пространства или кватернионных подмногообразиї кватернионного пространства, индуцированные римановыми связностями этих пространств.



Уравнение Янга - Миллса является вариационным уравнением Әйлера - Лагранжа для функционала $L(\Gamma)$ в пространстве связностей главного расслоения $\pi$, задаваемого формулой



\[

L(\nabla)=\int_{M}\left\langle R^{\nabla}, R^{\nabla}\right\rangle .

\]



Римавово многообразие $(M, g)$ предполагается компактным и ориентированным, а скобка $\langle.$, означает скалярное произведөние в слоях векторного расслоения $\mathfrak{q}(M) \otimes \Lambda^{2}, \quad$ определяемое $A d G$-инвариантным скалярным произведением в алгебре Ли $\mathfrak{g}$ группы $G$ и скалярным произведением в слоях расслоения $\Lambda^{2}$ 2-форм на $M$, индуцированным метрикой $g$. Таким образом, Я.М. П. суть критич. точки функции $L(\nabla)$. Я.- M. п.



\includegraphics[max width=\textwidth, center]{2023_07_09_08609cb348b493df6eb8g-1(2)}

ции $L$ в точке $\nabla$ положительно определев (и, следовательно, $\nabla$ есть локальныї минимум функции $L)$ и с л аб о у с т о й ч и в ы м, если 2-іі диффференциал неотрицательно определен. Известно, что не супествует слабо устоӥчивых Я.-М. ш., в произвольном нетривиальноч главном расслоении над стандартноӥ сфероӥ $S^{n}$ при $n \geqslant 5$. С другой стороны, при $n \geqslant 3$ риманова связность стандартної римановоӥ метрики факторпространства $S^{n / \Gamma}$ сферы по свободно действующеӥ нетривиальной конечної группе изометриї $\Gamma$ является устойчивым Я.-M.II. [5].



Наибольший интерес для физики представляет Я.-М.п. на четырехмерных римановых (а также лоренцевых) многообразиях. В этом случае оператор Ходжа : отобраякает пространство 2-форм на $M$ (со значениями в цроизвольном векторном расслоении) на себя, причем это отображение является инволютивным $\left(*^{2}=i d\right)$ и зависит только от ориентации и конформного класса метрики $g$. Связность $\nabla$ в главном расслоении над $M^{4}$ на3. с а м о д у а ль Б о й с в я $\mathbf{3}$ н о с. т ь ю илI и н с т а в т о о о (соответственно, а н т и с а мод у ально ї с вя зностью или антиинс т а н т о н о м), если ее 2-форма кривизны $R \nabla$ является собственным вектором олератора Ходжа с собственным значением 1 (соответственно, -1). В силу тождества Бъянки инстантоны и антиинстантоны являются Я.- М. І. Более того, они являются точками абсолютного минимума функции $L$. В случае главного расслоения над стандартной сферой со структурной группоӥі $S U(2), \quad S U(3)$ или $S U(4)$ все локальные минимумы функции $L$ исчерпываются инстантонами и антиинстантонами (и, следовательно, являются глобальными). Риманова связность риманова многообразия $M^{4}$ является инстантоном только для многообразий с группой голономии $G \subset \operatorname{Sp}(1)$. Все такие компактные многообразıя исчерпываются комплексными поверхностями К3.



Группа $G(\pi)=\Gamma\left(P_{*} \mathrm{dd} G\right)$ тождественных на базе автоморфизмов расслоения $\pi: P \rightarrow M$ наз. к а л и $6 \mathrm{p}$ ов о ч н о й г р у П П о й, Она өстөственным образом действует в множестве $C+(\pi)$ инстантонов расслоения $\pi$ с группой голономии $G$. ()акторпространство $M+(\pi)=$ $=C^{+}(\pi) / G(\pi)$ наз. П р о с т р а нс т в о м мод ул в ї неприводимых инстантонов расслоения л. Если структурная группа $G$ расслоөния л компактна и полупроста, а база $M^{4}$ является компактным ориентированным рімановым многообразием с неотрицательной ненулевої скалярной кривизной и самодуальным тензором конформной кривизны Вейля, то пространство модулеї $M^{+}(\pi)$ либо пусто, либо является многообразиям размерности



$\operatorname{dim} M+(\pi)=2 p_{1}(a(M))-1 / 2 \operatorname{dim} G\left(\chi\left(M^{4}\right)-\sigma\left(M^{4}\right)\right)$, где $p_{1}(\mathfrak{q}(M))$ - первое Понтряаина число расслоения $\mathfrak{g}(M), \quad$ а $\chi\left(M^{4}\right), \quad \sigma\left(M^{4}\right)$ - характеристика Эйлера Пуанкаре и сизнатура многообразия $M^{4}$.



Наиболее полные результаты получены в физически важном случае расслоений над стандартноӥ сферой $S^{4}$ с классическими компактными структурными группами $G$. В частности, получено описание всех инстантонов в таких расслоениях в чисто алгебрапч. терминах (напр., в терминах нек-рых модулей над грассмановой алгеброй или в терминах репений нек-рых кватернионных матричных уравнений, см. [11). Для случая группы $G=S p(1)$ известно явное описание всех инстантонов. Напр., Для $S p(1)$-расслоения $\pi_{1}: P_{1} \rightarrow S^{4}$ с Понтрягина числом 1 пространство модулей $M+\left(\pi_{1}\right) \approx \mathbb{R}+\times \mathbb{H}$, где $\mathbb{R}+$ - множество положительных чисөл, а $\mathbb{H}-$ множество кватернионов. Паре $(\lambda, a) \in R+\times \mathbb{H}$ отвечает инстантон, задаваемыӥ g-значноӥ 1-формоӥ связности



\[

A(x)=\operatorname{Im} \frac{\overline{u(x)} \cdot d u(x)}{1+|u(x)|^{2}},

\]



где $u(x)=\lambda(\overline{a-x})^{-1}, \quad x \in H$. Кватернионы из $\mathrm{H} \approx \mathbb{R}^{4}$ отождествляются с точками сферы $S^{4}$ с помощью стереографич. проекции, а алгебра ЈІи $\mathfrak{q}=S p(1)$ рассматривается как алгебра Ли Im Ч чисто мнимых кватернионов. Лит.: [1] М а І и І Ю. И., в сб.: Итоги науки и техники

Соврем. проблемы матем., т. 17, M., 1981, с. 3-55: [2]



\begin{center}

\includegraphics[max width=\textwidth]{2023_07_09_08609cb348b493df6eb8g-1(1)}

\end{center}



\begin{center}

\includegraphics[max width=\textwidth]{2023_07_09_08609cb348b493df6eb8g-1(4)}

\end{center}



\includegraphics[max width=\textwidth, center]{2023_07_09_08609cb348b493df6eb8g-1}

"Proc. Nat. Acad. Sci. USA", 1979 , v. 76, Ni 4, p. 1550-53; [6] Кон опл е в a H. П.. П о п о в В, Н,, Жалибровочные поля, 2 изд.. M., 1980; [7] Y a n g C. N. [8] M I ls R. L.. "идеи в физике, пер. с англ., м., 1983. Д. В. Алексеевский.



ЯЧЕЙКА - термин, применяемый иногда для обозначения такого параллелограмма периодов двоякопериодической бункции $f(z)$, на сторонах к-рого нет полюсов и к-рый получается из основного параллелограмма периодов сдвигом на нек-рыӥі вектор $z_{0} \in \mathbb{C}$.



\includegraphics[max width=\textwidth, center]{2023_07_09_08609cb348b493df6eb8g-1(3)}

современного анализа, пер. с англ., 2 изд., ч. 2, М., 1963.





\end{document}